\chapter{Einführung in die Medizin}

\emph{Maintainer: Sebastian Gabriel}

\minitoc

\gDefinition{Die Medizin ist die Lehre von Gesundheit und
Krankheit sowie die Kunde der Gesunderhaltung, Heilung und
Linderung.}

% \paragraph*{Einteilung} Notfälle kann man mehr oder weniger willkürlich
% einteilen
% 
% \begin{itemize}
% 	\item Störungen des Bewu{\ss}tseins --  neurologische Notfälle
% 	\item pulmonale Notfälle
% 	\item Kreislaufstörungen
% 	\item Gefä{\ss}verschlüsse
% 	\item Abdominelle Notfälle
% 	\item Urologische Erkrankungen
% \end{itemize}

\section{Begriffe}

Die Kenntnis von Fachbegriffen in der Medizin und
Sanitätshilfe ist nicht deshalb wichtig, weil die Fachwörter
so gut klingen. Vielmehr dient die Fachsprache dazu, dass
verschiedene Personen miteinander kommunizieren können und
dabei sichergestellt ist, dass sie sich gegenseitig
verstehen.

\paragraph{Symptome} \index{Symptom} Krankheitserscheinungen
\paragraph{Komplikationen} \index{Komplikation}
möglicherweise auftretende Schwierigkeiten
\paragraph{Therapie} \index{Therapie} Ma{\ss}nahmen zur
Heilung oder Linderung von Krankheiten und Symptomen
\paragraph{Diagnose} \index{Diagnose}

Der Begriff Diagnose wird dauernd verwendet, dennoch ist er
nicht einfach zu erklären. Eine Diagnose gibt einem Zustand
eine Bezeichnung, das heisst einem Beschwerdebild wird ein
definiertes Kranheitsbild zugeordnet.

\begin{tabularx}{\linewidth}{XcX}\toprule
Beschwerden und Symptome des Patienten
&
{\color{darkBlue}\sffamily\bfseries $\rightarrow$~DIAGNOSE~$\leftarrow$}
&
Definiertes Krankheitsbild
\\\addlinespace
Patient
&
{\color{darkBlue}\sffamily\bfseries $\rightarrow$~DIAGNOSE~$\leftarrow$}
&
Lehrbuch
\\\addlinespace
Realität
&
{\color{darkBlue}\sffamily\bfseries $\rightarrow$~DIAGNOSE~$\leftarrow$}
&
Lehre
\\\bottomrule
\end{tabularx}

\gFazit{Die Diagnose ordnet die Beschwerden des Patienten einem definierten
Krankheitsbild zu.}

Mitunter kann diese Zuordnung einfach sein, sehr oft ist die Diagnosestellung
aber unsicher. Um die Unsicherheit dieser Diagnosen
anzuzeigen haben sich verschiedene Begriffe eingebürgert: 

\begin{description}
    \item[Verdachtsdiagnose] \index{Verdachtsdiagnose} Die
    Diagnose ist unsicher, man hat zwar den Verdacht in Richtung eines bestimmten Krankheitsbildes, andere ähnliche
    Krankheitsbilder können aber nicht mit hinreichender Sicherheit
    ausgeschlossen werden. Die endgültige Klärung überlässt man jemand
    anderen\footnote{z.B. dem Krankenhaus}. \par 
    Sanitäter erstellen i.d.R. nur Verdachtsdiagnosen.
     
    \item[Arbeitsdiagnose] \index{Arbeitsdiagnose}
    Arbeitsdiagnosen sind Verdachtsdiagnosen bei denen eine endgültige Klärung nicht möglich ist. Für die Behandlung wird die
    wahrscheinlichste Diagnose ausgewählt um \emph{``am Patienten arbeiten zu
    können''.}
    \item[Notfalldiagnose] \index{Notfalldiagnose} Eine
    Notfalldiagnose ist im Prinzip eine Arbeitsdiagnose. Der Patient ist dabei in einem Zustand bei dem
    \gEs{sofortiges Handeln} notwendig ist und die eigentliche Diagnose, die zu
    dieser Situation geführt hat, in dem Moment unwichtig ist. \par
    Ein Besipiel hierfür sind die Bewußtlosigkeit und der Kreislaufstillstand.
    Die Ursachen dafür sind erstmal unwichtig, m
   it der entsprechenden Behandlung der Notfalldiagnose muß
    sofort begonnen werden.
    \item['Status post' -- 'Zustand nach']
    \index{Status post}\index{Zustand
    nach}\index{St.p.|see{Status post}}
    \index{Z.n.|see{Zustand nach}} Abk. \gE{'St.p.'}, \gE{'Z.n.'}. Bezeichnet eine ``alte Diagnose'', d.h. eine Erkrankung oder eine Behandlung die einmal durchgemacht wurde. Z.B.:
    \begin{itemize}
        \item \gTt{Z.n. Blinddarm-Operation}: Der Patient
        hatte einmal eine Blinddarm-Operation.
        \item \gTt{Z.n. Herzinfarkt 2/2006}: Der Patient
        hatte im Februar 2006 einen Herzinfarkt.
    \end{itemize}
    \emph{``Umgangssprachlich''} wird oft auch der Umstand der zu einer
    Diagnose geführt hat mit \emph{``Z.n.''} angegeben. 
    \begin{itemize}
        \item \gTt{V.a. Schenkelhalsfraktur re., Z.n. Sturz}: Der
        Patient Patient ist gestürzt und hat deshalb möglicherweise eine
        Schenkelhalsfraktur rechts.
    \end{itemize}
\end{description}
\paragraph{Differentialdiagnosen (DD)} \index{} Diagnosen,
welche bei den jeweiligen Symptomen auch
möglich, aber nicht so wahrscheinlich sind wie jene
Diagnose für welche man sich ``entscheidet''.

Beispiel: Patient mit Kopfschmerzen nach ausgiebigen
Alkoholgenuß am Vortag.
\begin{description}
    \item[Verdachtsdiagnose:] \gE{``Kater''}
    \item[Differentialdiagnosen:]
    Migraine, Hirnblutung, Schlaganfall,
    Hypoglykämie, Intoxikation, \ldots
\end{description}

% TODO Migraine? Migräne?

\paragraph{Diagnostik}\index{Diagnostik}

Als Diagnostik versteht man alle Bemühungen um eine Diagnose zu stellen, egal
ob man sich der eignen Sinne oder komplizierter Geräte
bedient.

\paragraph{Befund}\index{Befund}

Jede diagnostische Maßnahme hat (hoffentlich) ein Ergebnis, welches man als
\gT{Befund} bezeichnet. 

%\subsection{Medizinische Fachbereiche}

%\subsection{Andere Medizinische Berufe --- Ein kleines
%\emph{"`who is who"'}}

\section{\gZ{Reserviert\gA{??}}}

\section{\gZ{Reserviert\gA{??}}}






\section{Dokumentation -- Von der medizinischen Seite
betrachtet}\index{Dokumentation!medizinische Aspekte}

Ihre Dokumentation muss eine Geschichte erzählen. Eine
Geschichte die ein Dritter, unbeteiligter, verstehen
\gE{muss}. Dieser unbeteiligte Dritte soll allein anhand der
Schilderung in Ihrer Dokumentation fähig sein eine
Verdachtsdiagnose zu stellen. Weiters sollte er
nachvollziehen können, was mit dem Patient passiert ist.
Auch die Ereignisse davor oder die Einsatzsituation soll
für denjenigen verständlich sein (Unfallhergang,
Wohnsituation bei psychiatrischen Patienten, \ldots)

Schlecht:
\begin{quote}

{\sffamily Diagnose:} {\ttfamily RQW}

{\sffamily Befunde:} {\sffamily RR} {\ttfamily 130/70,}
{\sffamily HF} {\ttfamily 90} {\sffamily SpO2}
{\ttfamily 98\% }

{\sffamily Anmerkungen: --}
\end{quote}

Besser:

\begin{quote}
{\sffamily Diagnose:} {\ttfamily RQW Stirn}

{\sffamily Befunde:} {\sffamily RR} {\ttfamily 130/70,}
{\sffamily HF} {\ttfamily 90} {\sffamily SpO2}
{\ttfamily 98\% }

{\sffamily Anmerkungen:} {\ttfamily Pat. fuhr mit Fahrrad
und prallte mit Stirn gegen Ast., kein Sturz. Lt. Begleiter
keine B.losigkeit \underline{Traumacheck:} RQW Stirn,
ca 50\,ml Blutverlust sonst unauffällig. \underline{Neuro:}
grob neurologisch unauffällig, Bew. klar, Pat zeitl, örtl, zur Situation \& Person orientiert, Pupillen
mittelweit, Reakt. prompt und seitengleich, Kraft\&Gefühl
seitengl. in OE und UE. \underline{Anamn.:} \underline{Sympt:} Kopfschmerz seit
Aufprall. \underline{Krankh.:} keine bek. \underline{All.}
keine bek. \underline{Med.:} Aspirin heute früh. letzte
Tetanus-Impf. nicht bek. \underline{Ereignisse} s.o.
\underline{Letzte} Mahlzeit: heute Früh 2 Semmel\gA{ KOCH: ? muss sein? GABS:
ja, 1) wegen s-kamel, 2) wegen nüchtern } 

\underline{Maßnahmen:} Steriler Verband, nüchtern belassen,
Transport liegend mit erh OK.}
\end{quote}

Psychiatrisch:

\begin{quote}
{\sffamily Diagnose:} {\ttfamily Psychiatrische Erkankung,
V.a. Wahnvorstellungen, Unterbringung nach UbG}

{\sffamily Befunde:} {\sffamily RR} {\ttfamily 130/70,}
{\sffamily HF} {\ttfamily 90} {\sffamily SpO2}
{\ttfamily 98\% }

{\sffamily Anmerkungen:} {\ttfamily Pat. sitzend in Wohnung angetr., von Pol. mit Handschellen
gefesselt. Pat. warf lt. Polizei Gegstände aus d Fenster
auf Straße, verletzte Passanten ($\rightarrow$ FLO-1). Pat.
gibt an sich geg. Geheimdienst gewehrt zu haben,
jetzt agitiert und unkooperativ. 

\underline{Traumacheck:} keine äußeren Anz f Verletzungen,
sonst nicht erhhebbar (unkoop.). \underline{Neuro:} Pat.
ist wach, sonst nicht erhebbar

\underline{Anamn.:} \underline{Sympt:} keine Antwort
\underline{Krankh.:} lt. Spitalbrief Otto-Wagner-Sp. v.
1.4.08 paranoide Schizophrenie \underline{All.} k.A.
\underline{Med.:} nicht eruierbar, auf Tisch mehrere
angebr. Packungen Rohypnol \underline{Ereignisse} s.o.
\underline{Letzter Spitalsaufenthalt} k.A.

\underline{Wohnsituation:} heruntergekommen, vermüllt, Pat.
wohnt mit großem Hund, ganze Wohnung verkotet, bestialisch
stinkend. 

\underline{Maßnahmen:} Einweisung nach UbG, Transport in
Handschellen in Pol-begl. (Viktor 6, DNr. 4711). Protokoll
Amtsarzt ad. Spital.}
\end{quote}

\section{\gZ{Reserviert\gA{Medizinische Fachbereiche}}}

\section{\gZ{Reserviert\gA{Medizinische Berufsgruppen}}}

\section{\gZ{Reserviert}}